\section{Related Work}

\subsection{Vision-Language-Action Models}

Vision-language-action models (VLAs) adapt pretrained vision-language models for closed-loop robot control, with the aim of transferring semantic knowledge acquired from large-scale vision-language pretraining into manipulation policies. RT-2 established an influential formulation in which robot actions are represented as text tokens and the model is co-fine-tuned on vision-language and robot-trajectory data. Its evaluations showed improved generalization to novel objects and instructions, including commands involving relative size and spatial proximity \cite{brohan2023rt2}.

More recent VLAs explore different interfaces between pretrained representations and robot actions. $\pi_{0.5}$ extends $\pi_0$ through heterogeneous co-training on robot data, high-level semantic prediction, verbal instructions, and web data \cite{black2025pi05}. MolmoAct2 connects an embodied-reasoning vision-language model to a flow-matching continuous-action expert through per-layer conditioning \cite{fang2026molmoact2}. VLA-0 instead preserves the base vision-language-model architecture and expresses continuous actions as numerical text strings \cite{goyal2025vla0}. These systems differ not only in their action representations, but also in their pretraining data and objectives. We therefore compare complete pretrained VLA families after adaptation and do not interpret differences between them as causal evidence about any single architectural choice.

Collectively, these models motivate evaluating whether semantic knowledge remains usable after adaptation to robot control. However, successful task execution alone does not establish that a policy grounded the instruction in the current scene: the same outcome may result from recurring object identities, layouts, target positions, or familiar action templates. Spa-Bench targets this distinction by varying the task-relevant semantic variable while retaining familiar objects and manipulation behaviors.

\subsection{Reasoning and Generalization Benchmarks}

Existing benchmarks probe language-conditioned manipulation at different scales and under different forms of distribution shift. VLABench emphasizes breadth, spanning spatial understanding, semantic instruction following, physical-law understanding, knowledge transfer, and long-horizon reasoning across a large collection of tasks and objects \cite{zhang2024vlabench}. Spa-Bench is complementary to this broad evaluation: it uses a smaller physical workspace and a restricted set of manipulation skills to isolate concept--argument transfer. In particular, it tests whether a model can recombine individually familiar concepts, objects, and physical behaviors when their exact combination was withheld during fine-tuning---an axis that broad task-level success rates do not isolate directly.

LIBERO-PRO focuses on evaluation reliability and shows that performance under standard protocols can deteriorate after controlled changes to manipulated objects, initial states, instructions, or environments \cite{zhou2025liberopro}. This finding motivates our matched and counterfactual scene families, which vary the instruction or task-relevant scene variable while keeping manipulation demands as comparable as possible. Matched manipulation controls further test whether a primary-transfer failure is better explained by object recognition or physical execution than by semantic selection.

STAR-Gen provides a taxonomy that distinguishes changes to visual input, semantic instruction, and required behavior \cite{gao2026stargen}. Under this taxonomy, Spa-Bench primarily studies semantic and semantic-plus-behavioral generalization. Its cross-task-object condition introduces a controlled visual shift by assigning a familiar concept to an object that did not serve that role during task-specific training, whereas its language-robustness trials change the wording while holding the intended behavior and scene fixed. Because aggregate scores can obscure differences among these axes, we report concept/compositional transfer, language robustness, cross-task-object performance, and no-intervention behavior separately rather than pooling all 900 rollouts into a single measure.

Spa-Bench is not intended as a benchmark of long-horizon planning, broad open-world competence, or vision-language question answering without physical execution. Instead, it is a controlled real-robot diagnostic of whether an end-to-end policy binds instruction-level semantic variables to task-relevant scene variables and expresses that binding through action. Spa-Bench therefore complements existing Spa-Benchs by isolating scene-grounded semantic transfer under controlled training exclusions, rather than measuring only broad task success, environmental robustness, or long-horizon planning ability.
